\documentclass[11pt,a4paper]{article}
\usepackage[T2A]{fontenc}
\usepackage[utf8]{inputenc}
\usepackage[russian]{babel}
\usepackage{amsmath,amssymb,amsthm,mathtools}
\usepackage[margin=2.3cm]{geometry}
\usepackage{enumitem}
\usepackage{tikz}
\usepackage[colorlinks=true,linkcolor=blue!50!black,citecolor=blue!50!black,urlcolor=blue!50!black]{hyperref}

\newtheorem{theorem}{Теорема}[section]
\newtheorem{lemma}[theorem]{Лемма}
\newtheorem{proposition}[theorem]{Предложение}
\newtheorem{corollary}[theorem]{Следствие}
\theoremstyle{definition}
\newtheorem{definition}[theorem]{Определение}
\newtheorem{example}[theorem]{Пример}
\newtheorem{remark}[theorem]{Замечание}

\newcommand{\R}{\mathbb R}
\newcommand{\eps}{\varepsilon}
\newcommand{\sym}{\operatorname{sym}}
\newcommand{\Id}{\mathrm{Id}}
\newcommand{\mk}[1]{\textsf{[#1]}}
\newcommand{\term}[1]{\emph{#1}}

\title{Неустойчивость равновесия в аналитическом потенциальном поле на плоскости\\[4pt]
\large Обзор работы В.\,П.~Паламодова (1977) для читателя вне теории устойчивости}
\author{Главный агент исследования (Claude) по поручению владельца}
\date{21.10.2026, четвёртая редакция}

\begin{document}
\maketitle

\begin{abstract}
Пусть $U$ --- вещественно-аналитическая функция на плоскости, $0$ --- её критическая точка и не
точка минимума. Тогда положение равновесия $x=0$ системы $\ddot x=-\nabla U(x)$ неустойчиво.
Это доказал В.\,П.~Паламодов в 1977 году (независимо --- С.~Талиаферро, 1980). Текст объясняет,
что это значит и как доказывается: какие понятия механики здесь участвуют, откуда берётся
неустойчивость, где работает аналитичность и почему тот же путь пока не пройден в размерности
три и выше. Специальные термины определены по ходу; их список --- в приложении~\ref{app:glossary}.
\end{abstract}

\paragraph{Состояние текста.} Разделы \ref{sec:mech}--\ref{sec:method} и \ref{sec:after} проверены
независимым рецензентом (отзыв: \texttt{notes/reviewer/2026-10-21-palamodov1977-obzor.md}), эта
редакция учитывает его замечания. Раздел~\ref{sec:plane} написан по тексту статьи [P77]
и проверен рецензентом (отзыв: \texttt{notes/reviewer/2026-10-21-palamodov1977-obzor-v3.md}); эта
редакция учитывает его замечания.

\paragraph{Метки.} \mk{Литература: ссылка} --- опубликованный результат; если первоисточник не
открывался, это оговорено. \mk{Доказано} --- полное доказательство приведено здесь.
\mk{Эскиз} --- идея с перечисленными пробелами. Доказательства этой редакции после правок
ещё ждут повторной проверки.

\tableofcontents

%=====================================================================
\section{Механика: движение, энергия, устойчивость}\label{sec:mech}

\subsection{Уравнение движения и фазовое пространство}

Точка (частица) единичной массы движется в $\R^n$. Её \term{положение} $x\in\R^n$, \term{скорость}
$\dot x=dx/dt$, ускорение $\ddot x$. На неё действует сила, заданная \term{потенциалом} (иначе:
\term{потенциальной энергией}) --- функцией $U\colon W\to\R$ на открытом множестве $W\subset\R^n$:
сила равна $-\nabla U(x)$. Закон Ньютона даёт
\begin{equation}\label{eq:newton}
\ddot x=-\nabla U(x).
\end{equation}
Пространство положений ($W$) называют \term{конфигурационным пространством}, $n$ --- \term{числом
степеней свободы}. Уравнение \eqref{eq:newton} равносильно системе первого порядка
\begin{equation}\label{eq:system}
\dot x=p,\qquad \dot p=-\nabla U(x)
\end{equation}
на \term{фазовом пространстве} $W\times\R^n$; его точки --- пары (положение $x$, скорость $p$)
(при единичной массе скорость совпадает с импульсом, отсюда буква $p$). \term{Движение} ---
решение $x(t)$ уравнения \eqref{eq:newton}; \term{фазовая траектория} --- кривая
$t\mapsto(x(t),\dot x(t))$. Если $U\in C^2$, правая часть \eqref{eq:system} локально липшицева, и
через каждую точку фазового пространства проходит ровно одно решение (теорема Пикара); оно
продолжается, пока не покидает компактные подмножества $W\times\R^n$.

Такие системы называют \term{натуральными механическими системами}: энергия складывается из
кинетической $\tfrac12|\dot x|^2$ и потенциальной $U(x)$. В общей постановке кинетическая
энергия --- положительно определённая квадратичная форма скорости, зависящая от положения;
у нас она евклидова. О знаках: в части литературы используется \term{силовая функция} $-U$, и
тогда <<минимум потенциала>> превращается в <<максимум силовой функции>>.

\subsection{Закон сохранения энергии}

\begin{lemma}\label{lem:energy}
\mk{Доказано} Функция $H(x,p)=\tfrac12|p|^2+U(x)$ (\term{полная энергия}) постоянна вдоль
каждого решения \eqref{eq:system}.
\end{lemma}
\begin{proof}
$\frac{d}{dt}H(x,\dot x)=\langle\dot x,\ddot x\rangle+\langle\nabla U(x),\dot x\rangle
=\langle\dot x,-\nabla U(x)+\nabla U(x)\rangle=0$.
\end{proof}

Значение $E=H$ на движении называют его \term{энергией}. Из $\tfrac12|\dot x|^2=E-U(x)\ge0$
следует, что движение с энергией $E$ не выходит из множества $\{x:U(x)\le E\}$ (\term{область
возможного движения}), а модуль скорости определяется положением: $|\dot x|=\sqrt{2(E-U(x))}$.
Это главный инструмент всего текста: при $E<0$ движение целиком проходит там, где $U\le E<0$.

\subsection{Положение равновесия и устойчивость по Ляпунову}

\term{Положение равновесия} --- постоянное движение $x(t)\equiv x_0$. По \eqref{eq:newton} это
возможно ровно тогда, когда $\nabla U(x_0)=0$, т.\,е. $x_0$ --- \term{критическая точка}
потенциала. Без ограничения общности $x_0=0$ и $U(0)=0$. Критическая точка \term{изолирована},
если в некоторой проколотой окрестности нуля $\nabla U\ne0$.

\begin{definition}[устойчивость по Ляпунову]\label{def:stab}
Равновесие $(x,p)=(0,0)$ \term{устойчиво по Ляпунову}, если для любого $\eps>0$ найдётся
$\delta>0$ такое, что каждое решение с $|x(0)|+|\dot x(0)|<\delta$ определено при всех $t\ge0$
и удовлетворяет $|x(t)|+|\dot x(t)|<\eps$ при всех $t\ge0$. Иначе равновесие
\term{неустойчиво}: существует $\eps>0$ (его можно считать настолько малым, что шар
$\{|x|\le\eps\}$ лежит в $W$) такое, что при любом $\delta>0$ найдётся решение с
$|x(0)|+|\dot x(0)|<\delta$, которое в какой-то момент $t>0$ достигает $|x(t)|+|\dot x(t)|\ge\eps$.
\end{definition}

Иначе говоря, устойчивость означает: сколь угодно малое возмущение равновесия (по положению и
по скорости) остаётся малым навсегда. Неустойчивость: есть фиксированный масштаб $\eps$, до
которого раскачиваются некоторые сколь угодно малые возмущения.

\begin{example}
$n=1$. При $U=x^2/2$ (пружина) решения $x=a\cos t+b\sin t$ --- равновесие устойчиво.
При $U=-x^2/2$ (перевёрнутый маятник) решение $x=a\cosh t$ уходит от нуля --- неустойчиво.
\end{example}

Замена $t\mapsto-t$, $p\mapsto-p$ переводит решения в решения (\term{обратимость}), поэтому всё
сказанное о будущем верно и о прошлом.

\subsection{Теорема Лагранжа--Дирихле}

\begin{theorem}[Лагранж 1788, Дирихле 1846]\label{thm:LD}
\mk{Литература: классический результат; доказательство стандартное и приведено полностью ---
\mk{Доказано}} Если $0$ --- точка \term{строгого локального минимума} $U$ (т.\,е. $U(x)>U(0)=0$
при $0<|x|\le r$), то равновесие устойчиво по Ляпунову.
\end{theorem}
\begin{proof}
Пусть $0<\eps\le r$ и $m(\eps)=\min_{|x|=\eps}U>0$. По непрерывности $H$ найдётся $\delta>0$
такое, что из $|x|+|p|<\delta$ следуют $|x|<\eps$ и $H(x,p)<m(\eps)$. Решение с такими
начальными данными, пока определено, не может достичь сферы $|x|=\eps$: там было бы
$H\ge U\ge m(\eps)$. Значит, $|x(t)|<\eps$, а так как $U\ge0$ в шаре, $|\dot x(t)|^2=2(H-U)<2m(\eps)$.
Фазовая точка остаётся в компакте, поэтому решение продолжается на все $t\ge0$, и при всех $t$
выполнено $|x|+|\dot x|<\eps+\sqrt{2m(\eps)}$. Эта величина стремится к нулю при $\eps\to0$.
\end{proof}

Механический смысл: в ямке потенциала частица с малой энергией заперта стенками ямки.

\subsection{Обратная задача и задача Арнольда}

Вопрос, восходящий к Ляпунову: верно ли обратное --- если $0$ не точка минимума $U$, то
равновесие неустойчиво? Различают два случая.

\begin{itemize}[leftmargin=*]
\item \term{Нестрогий минимум}: $U\ge0$ около нуля, но нули $U$ подходят к $0$. Пример:
$U=x_1^2$ на $\R^2$; тогда $x_1\equiv0$, $x_2=ct$ --- решение, и равновесие неустойчиво.
\item \term{Не минимум}: в любой окрестности нуля есть точки, где $U<0$.
\end{itemize}

Если критическая точка $0$ изолирована, нестрогий минимум невозможен: нуль $q\ne0$ функции
$U\ge0$ сам был бы точкой минимума, а значит критической точкой. \mk{Доказано}

\paragraph{Задача Арнольда 1971-4, случай (A).} В сборнике <<Задачи Арнольда>> задача 1971-4
цитируется так \mk{Литература: дословно в arXiv:1906.11962, 2208.01139, 2108.05829; сама книга не
открывалась}: <<Prove the instability of the equilibrium 0 of the analytic system
$\ddot x=-\partial U/\partial x$ in the case where the isolated critical point 0 of the potential U
is not a minimum>>. То есть:

\medskip
\noindent\fbox{\parbox{0.96\linewidth}{\textbf{(A)} $U$ вещественно-аналитична около
$0\in\R^n$, $0$ --- изолированная критическая точка, не минимум. Доказать, что равновесие
неустойчиво по Ляпунову.}}
\medskip

Для $n=1$ это просто (предложение~\ref{prop:n1}), для $n=2$ доказано Паламодовым [P77] и
Талиаферро [Ta80b]; при $n\ge3$ задача открыта \mk{Литература: Burgos--Paternain,
arXiv:2208.01139, введение; Burgos, J.~Differential Equations 454 (2026), arXiv:2506.19135, §1.2}.

\subsection{Почему важна аналитичность}\label{subsec:smooth}

\begin{example}[Пенлеве, Уинтнер]\label{ex:PW}
$n=1$, $U(x)=e^{-1/x^2}\cos(1/x)$, $U(0)=0$. \mk{Литература: Painlevé 1904, Wintner 1941; по
вторичным источникам. Проверка ниже --- \mk{Доказано}} Эта функция бесконечно
дифференцируема, но не аналитична в нуле: все её производные в нуле равны нулю, а сама она не
тождественный нуль. Точка $0$ не минимум: в точках $x=1/((2k+1)\pi)$ значение отрицательно. Но
равновесие устойчиво. В точках $\pm a_k$, $a_k=1/(2k\pi)$, значения $U(\pm a_k)=e^{-4k^2\pi^2}>0$ ---
это <<барьеры>>. Частица на прямой с энергией $E<U(\pm a_k)$, стартовавшая в $(-a_k,a_k)$, не
может пройти через $\pm a_k$: там было бы $\tfrac12\dot x^2=E-U<0$. Если начальные данные малы,
то мало и $E$, и частица заперта в сколь угодно малом интервале $(-a_k,a_k)$ со скоростью
$|\dot x|\le\sqrt{2(E+\max|U|)}$. Критическая точка $0$ здесь \emph{не} изолирована: между
соседними нулями $U$ лежит критическая точка (теорема Ролля), и они накапливаются к нулю.
\end{example}

\begin{proposition}\label{prop:n1}
\mk{Доказано} Случай (A) при $n=1$ верен даже для $U\in C^1$.
\end{proposition}
\begin{proof}
Так как $U'\ne0$ на $(0,r)$ и на $(-r,0)$, $U$ строго монотонна на каждом из этих интервалов.
Раз $0$ не минимум, на одном из них, скажем на $(0,r)$, $U$ убывает и отрицательна. Частица,
отпущенная с нулевой скоростью из $x_0\in(0,r/2)$, имеет $\ddot x=-U'(x)>0$, поэтому движется
вправо, и $\tfrac12\dot x^2=U(x_0)-U(x)$ растёт. Остановиться или стремиться к точке внутри
$(0,r)$ она не может: предельная точка была бы положением равновесия. Значит, она достигает
$x=r/2$, сколь бы мало ни было $x_0$.
\end{proof}

При $n\ge2$ гладкие примеры тоньше. Лалуа (1976) построил гладкий $U$ на плоскости без минимума
в нуле с устойчивым равновесием; критическая точка там тоже не изолирована \mk{Литература:
Laloy, Int.\ J.\ Non-Linear Mech.\ 11 (1976); формула и свойства --- по вторичным источникам}.

Даже неустойчивость \emph{максимума} гладкого потенциала не очевидна: энергия ничего не
запрещает, и нужно показать, что какие-то движения действительно уходят. Её доказали Хагедорн
(1971) и Талиаферро (1980) \mk{Литература: Hagedorn, ARMA 42 (1971); Taliaferro, ARMA 73 (1980);
по вторичным источникам}. Уренья (2020) построил гладкий $U$ на плоскости со строгим глобальным
максимумом в нуле, $\nabla U\ne0$ вне нуля, и с последовательностью нетривиальных периодических
движений, стягивающихся к равновесию \mk{Литература: Ureña, Ann.\ Mat.\ Pura Appl.\ 199 (2020),
теорема 1.1 в цитате arXiv:2208.01139; в статье речь о неустойчивости в смысле Зигеля--Мозера}.
Burgos и Paternain делают отсюда вывод, что в гладком классе неверна и тотальная неустойчивость
(определение~\ref{def:total}) \mk{Литература: arXiv:2208.01139, введение}; для этого нужно, чтобы
энергии этих орбит были ниже $U(0)$, а в формулировке теоремы Уреньи об энергиях не сказано ---
так что этот вывод мы считаем вторичным.

Вывод: в доказательстве (A) аналитичность (или изолированность) обязана где-то работать по
существу. Есть ли гладкий пример с \emph{изолированной} критической точкой без минимума и с
устойчивым равновесием, неизвестно \mk{Литература: в просмотренной литературе нет ни примера, ни
утверждения о его отсутствии; см. \texttt{problem/KNOWN\_RESULTS.md}, §3}.

\subsection{Три способа доказать неустойчивость}\label{subsec:three}

\paragraph{1. Асимптотическая траектория.} Так называют движение $x(t)\not\equiv0$, для
которого $(x(t),\dot x(t))\to(0,0)$ при $t\to-\infty$ (или при $t\to+\infty$).

\begin{lemma}\label{lem:asympt}
\mk{Доказано} Если есть асимптотическая траектория, равновесие неустойчиво.
\end{lemma}
\begin{proof}
По обратимости можно считать, что $(x,\dot x)\to0$ при $t\to-\infty$. Положим
$\eps=|x(0)|+|\dot x(0)|>0$. Для любого $\delta>0$ найдётся $t_0<0$ с
$|x(t_0)|+|\dot x(t_0)|<\delta$; решение, стартующее в момент $t_0$, к моменту $0$ достигает
$\eps$.
\end{proof}

Этим путём шли Козлов, Талиаферро и другие: строили такое решение в виде рядов по
отрицательным (вообще говоря, дробным) степеням $t$ \mk{Литература: см.
\texttt{problem/KNOWN\_RESULTS.md}, §2; по вторичным источникам}.

\paragraph{2. Функция Четаева.}
\begin{theorem}[Четаев]\label{thm:chetaev}
\mk{Литература: Н.\,Г.~Четаев, классическая теорема о неустойчивости; формулировка стандартная}
Пусть $V$ --- $C^1$-функция в окрестности $\Omega$ равновесия в фазовом пространстве, $D$ ---
связная компонента множества $\{V>0\}\cap\Omega$, и равновесие лежит на границе $D$. Пусть $V$
ограничена на $D$, а её производная вдоль решений $\dot V$ положительна на $D$ и отделена от нуля
на каждом множестве $D\cap\{V\ge\alpha\}$, $\alpha>0$. Тогда равновесие неустойчиво.
\end{theorem}
Идея: решение, стартовавшее в $D$ сколь угодно близко к равновесию, не может покинуть $D$
через часть границы, где $V=0$ ($V$ растёт), и не может оставаться в $\Omega$ вечно ($V$ росла бы
неограниченно). Значит, оно выходит на границу $\Omega$. Метод Паламодова --- частный случай этой
идеи с функцией $V=\langle\dot x,v(x)\rangle$ (раздел~\ref{sec:method}).

\paragraph{3. Тотальная неустойчивость.}

\begin{definition}[тотальная неустойчивость]\label{def:total}
Критическая точка $0$ \term{тотально неустойчива}, если существуют окрестность $W_0$ нуля в
конфигурационном пространстве и функция $T\colon(-\infty,0)\to(0,\infty)$ такие, что каждое
движение с энергией $E<0$ проводит в $W_0$ без перерыва не больше времени $T(E)$: если
$x(t)\in W_0$ при $t\in[t_1,t_2]$, то $t_2-t_1\le T(E)$.
\end{definition}

\begin{remark}
\mk{Литература} В такой форме (<<с функцией $T(E)$>>) утверждение анонсировал Паламодов в 1995
году для всех $n$; так его цитирует Burgos (arXiv:2108.05829v4, введение). В обзоре Паламодова
2020 года [P20] определение формулируется иначе (по пересказу; первоисточник не сверен), у
Burgos (2025) --- иначе и слабее. Нам нужна форма определения~\ref{def:total}: её даёт метод.
\end{remark}

\begin{lemma}\label{lem:total-lyap}
\mk{Доказано} Если $0$ тотально неустойчива и не минимум, равновесие неустойчиво по Ляпунову.
Изолированность для этого не нужна.
\end{lemma}
\begin{proof}
Пусть шар $\{|x|<\eps\}$ лежит в $W_0$. Возьмём $x_0$ с $|x_0|<\delta$ и $U(x_0)<0$ (есть, так
как $0$ не минимум) и нулевую начальную скорость. Энергия движения $E=U(x_0)<0$, поэтому за
время $T(E)$ оно покидает $W_0$, а значит, достигает $|x|\ge\eps$.
\end{proof}

Тотальная неустойчивость сильнее, чем нужно для (A): она говорит обо \emph{всех} движениях с
отрицательной энергией, а не об одном. Зато её удобно доказывать <<одной функцией>>, и именно её
дают метод Паламодова и все результаты раздела~\ref{sec:method}. \emph{Гипотеза Паламодова}:
всякая критическая точка аналитического потенциала тотально неустойчива \mk{Литература: [P20],
гипотеза 1, по пересказу; не сверено}.

%=====================================================================
\section{Что нужно знать об аналитических функциях}\label{sec:analytic}

\subsection{Ряд Тейлора и младшая форма}

Функция $U$ \term{вещественно-аналитична} около $0$, если в некотором шаре она равна сумме
сходящегося степенного ряда. Группируя члены по степеням,
\[
U=u_m+u_{m+1}+u_{m+2}+\dots,
\]
где $u_k$ --- \term{однородная форма} степени $k$ (однородный многочлен: $u_k(\lambda x)=
\lambda^ku_k(x)$), а $u_m\not\equiv0$ --- \term{младшая форма}; $m$ --- \term{порядок} $U$ в нуле.
Если $0$ --- критическая точка и $U(0)=0$, то $m\ge2$. Форма $u_2$ задаётся \term{гессианом}
(матрицей вторых производных) $D^2U(0)$: $u_2(x)=\tfrac12\langle D^2U(0)x,x\rangle$.

\begin{itemize}[leftmargin=*]
\item Если $u_m$ принимает отрицательные значения, то и $U$ принимает: $U(tx)=t^m(u_m(x)+O(t))$.
\item Если $u_m>0$ вне нуля, то $0$ --- строгий минимум.
\item Трудный случай: $u_m\ge0$, но обращается в нуль на некоторых прямых. Тогда младшая форма
не решает дела. Пример: $U=y^2-x^4$ ($u_2=y^2$) --- не минимум; $U=y^2+x^4$ --- минимум.
\end{itemize}

Почти все классические критерии неустойчивости (Ляпунов --- по гессиану; Четаев --- для однородных
$U$; Козлов--Паламодов --- когда младшая форма принимает отрицательные значения) работают, когда
отсутствие минимума видно по младшей форме или по другой <<главной части>> ряда
\mk{Литература: см. \texttt{problem/KNOWN\_RESULTS.md}, §2; по вторичным источникам}. Ценность
результата Паламодова для $n=2$ в том, что на плоскости он не требует \emph{никаких} условий
такого рода.

\subsection{Полуаналитические множества и критические значения}

\term{Полуаналитическое множество} в окрестности точки --- конечное объединение множеств вида
$\{f_1=\dots=f_k=0,\ g_1>0,\dots,g_l>0\}$ с аналитическими $f_i,g_j$. Пример:
$\{U<0\}$, $\{\nabla U=0\}$.

\begin{theorem}[лемма о выборе кривой]\label{thm:curve}
\mk{Литература: J.~Milnor, Singular points of complex hypersurfaces, 1968, лемма 3.1 (для
полуалгебраических множеств); для полуаналитических --- Брюа--Картан, Лоясевич; не сверено с
первоисточником} Если точка $0$ лежит в замыкании полуаналитического множества $S$, то
существует вещественно-аналитическая (в том числе в точке $0$) кривая $\gamma\colon[0,\eps)\to\R^n$ с
$\gamma(0)=0$ и $\gamma(s)\in S$ при $s>0$.
\end{theorem}

\begin{lemma}\label{lem:critval}
\mk{Доказано, по модулю теоремы~\ref{thm:curve}} Пусть $U$ вещественно-аналитична около $0\in\R^n$
и $U(0)=0$. Тогда найдётся $\rho>0$ такое, что во всех критических точках $x\in B_\rho$
выполнено $U(x)=0$. Следовательно, для любого $E<0$
\[
\delta(E)=\min\{|\nabla U(x)|:\ |x|\le\rho/2,\ U(x)\le E\}>0
\]
(минимум по пустому множеству считаем равным $+\infty$).
\end{lemma}
\begin{proof}
Если бы критические точки $x_k\to0$ имели $U(x_k)\ne0$, то $0$ лежал бы в замыкании
полуаналитического множества $\{\nabla U=0\}\cap\{U\ne0\}$. По теореме~\ref{thm:curve} нашлась бы
кривая $\gamma$ из этого множества с $\gamma(0)=0$. Вдоль неё
$\frac{d}{ds}U(\gamma(s))=\langle\nabla U(\gamma(s)),\gamma'(s)\rangle=0$, так что
$U(\gamma(s))=U(0)=0$ --- противоречие. Второе утверждение: множество
$\{|x|\le\rho/2,\ U\le E\}$ компактно и не содержит критических точек (в его точках $U\le E<0$,
а критические значения в $B_\rho$ равны нулю), а $|\nabla U|$ непрерывна.
\end{proof}

Это первое место, где работает аналитичность. В примере Пенлеве--Уинтнера критические точки с
отрицательными значениями $U$ накапливаются к нулю, и лемма неверна.

\subsection{Нули и знаки аналитической функции на плоскости}\label{subsec:puiseux}

На плоскости устройство множества $\{U<0\}$ около нуля описывается полностью.

\begin{theorem}[Ньютон--Пюизё; локальная структура вещественной аналитической кривой]\label{thm:puiseux}
\mk{Литература: стандартный факт; см., например, C.\,T.\,C.~Wall, Singular points of plane curves,
2004, гл.~2; Brieskorn--Knörrer, Plane algebraic curves, §8.3; не сверено с первоисточником
в той форме, в какой оно здесь выписано} Пусть $U$ вещественно-аналитична около $0\in\R^2$,
$U(0)=0$, $U\not\equiv0$. Тогда при малом $\rho>0$ множество $\{U=0\}\cap B_\rho\setminus\{0\}$
--- объединение конечного числа $k\ge0$ попарно непересекающихся \term{полуветвей}
$\Gamma_1,\dots,\Gamma_k$. Каждая полуветвь пересекает каждую окружность $|x|=r$, $0<r<\rho$,
ровно в одной точке и в подходящих (повёрнутых) координатах задаётся как
\[
\Gamma_j=\{(x,y):\ 0<x<\rho_j,\ y=\textstyle\sum_{i\ge1}a_{ji}x^{i/N}\},
\]
где ряд по дробным степеням $x^{1/N}$ с \emph{вещественными} коэффициентами (\term{ряд Пюизё})
сходится, а поворот подобран так, что полуветвь касается в нуле положительной полуоси $x$ (в ряде
нет членов степени $\le1$).
\end{theorem}

Отсюда: проколотый круг $B_\rho\setminus\{0\}$ разрезан полуветвями на $k$ \term{криволинейных
секторов} (при $k=0$ --- он целиком один), и на каждом секторе $U$ сохраняет знак. Секторы, где
$U<0$, будем называть \term{рогами}. Если $0$ не минимум, есть хотя бы один рог. Рог может иметь
ненулевой угол раствора в нуле (например, квадрант у седла $U=xy$). Если ограничивающие сектор
полуветви касаются одного и того же луча, сектор либо узкий (<<клюв>>: ширина на расстоянии $r$
от нуля порядка $r^a$ с $a>1$), либо дополнение к такому клюву (например, у $U=x^3-y^2$ рог ---
всё, кроме клюва $\{|y|<x^{3/2}\}$).

\begin{example}\label{ex:horns}
\begin{enumerate}[leftmargin=*]
\item $U=-(x^2+y^2)$ (строгий максимум): $k=0$, весь проколотый круг --- один <<рог>>.
\item $U=xy$ (седло): $k=4$ (полуоси), два рога --- второй и четвёртый квадранты.
\item $U=y^2-x^4$: полуветви $y=\pm x^2$ при $x>0$ и при $x<0$; два рога $\{|y|<x^2\}$
вдоль оси $x$, ширина $\sim r^2$.
\item $U=(y-x^2)^2-x^6$: рога $\{|y-x^2|<|x|^3\}$ --- узкие ($\sim r^3$) и \emph{искривлённые}:
они идут вдоль параболы $y=x^2$. Нуль --- изолированная критическая точка
($\partial_yU=0\Rightarrow y=x^2\Rightarrow\partial_xU=-6x^5=0$).
\end{enumerate}
\end{example}

\begin{figure}[ht]
\centering
\begin{tikzpicture}[scale=2.6]
\draw[->,gray] (-0.1,0)--(1.6,0) node[right]{$x$};
\draw[->,gray] (0,-0.1)--(0,1.9) node[above]{$y$};
\fill[blue!15] plot[domain=0:1.35,samples=60,smooth] (\x,{\x*\x+0.2*\x*\x*\x}) -- plot[domain=1.35:0,samples=60,smooth] (\x,{\x*\x-0.2*\x*\x*\x}) -- cycle;
\draw[blue] plot[domain=0:1.35,samples=60,smooth] (\x,{\x*\x+0.2*\x*\x*\x});
\draw[blue] plot[domain=0:1.35,samples=60,smooth] (\x,{\x*\x-0.2*\x*\x*\x});
\draw[dashed,gray] plot[domain=0:1.32,samples=40,smooth] (\x,{\x*\x});
\fill[red] (1,0.9) circle (0.02);
\draw[->,red,thick] (1,0.9) -- (1.45,1.305);
\draw[dotted,red] (0,0) -- (1,0.9);
\node[red,right] at (1.45,1.25) {\small радиальное};
\node[red,right] at (1.45,1.12) {\small направление};
\node at (0.35,1.0) {\small $U>0$};
\node at (1.15,0.45) {\small $U>0$};
\draw[blue!60!black] (0.62,0.39) -- (0.95,0.2) node[right]{\small рог $U<0$};
\end{tikzpicture}
\caption{Искривлённый рог вдоль параболы $y=x^2$ (пунктир), как у функции $U=(y-x^2)^2-x^6$;
ширина рога для наглядности преувеличена. Из точки нижней половины рога радиальное направление
(красная стрелка, продолжение луча из нуля) идёт положе параболы и выходит из рога через нижнюю
границу, туда, где $U>0$.}
\label{fig:horn}
\end{figure}

%=====================================================================
\section{Метод Паламодова: функция $\langle\dot x,v(x)\rangle$}\label{sec:method}

\subsection{Образец: однородный потенциал и момент инерции}

Пусть $U$ однороден степени $m\ge2$ ($U(\lambda x)=\lambda^mU(x)$) и где-то отрицателен.
Рассмотрим $J(t)=\tfrac12|x(t)|^2$ --- половину \term{момента инерции} частицы относительно нуля.
По тождеству Эйлера $\langle x,\nabla U(x)\rangle=mU(x)$, поэтому
\[
\ddot J=|\dot x|^2+\langle x,\ddot x\rangle=|\dot x|^2-mU=2(E-U)-mU=2E-(m+2)U
\]
(\term{тождество Лагранжа--Якоби}; член $\langle x,\nabla U\rangle$ называют \term{вириалом} сил,
отсюда название <<вириальные>> у критериев этого типа). Если $E<0$, то $U\le E$ на движении, и
\[
\ddot J\ \ge\ 2E-(m+2)E=m|E|>0.
\]

\begin{proposition}\label{prop:homog}
\mk{Доказано} Для однородного $U$, принимающего отрицательные значения, равновесие тотально
неустойчиво и неустойчиво по Ляпунову.
\end{proposition}
\begin{proof}
Пусть движение с энергией $E<0$ проводит в шаре $|x|<\rho$ отрезок времени $[0,T]$. Тогда
$|\dot J|=|\langle x,\dot x\rangle|\le\rho\,\beta$, где $\beta=\sqrt{2(E-\min_{\bar B_\rho}U)}$ ---
оценка модуля скорости, и из $\ddot J\ge m|E|$ получаем $m|E|T\le\dot J(T)-\dot J(0)\le2\rho\beta$,
т.\,е. $T\le2\rho\beta/(m|E|)$. Это тотальная неустойчивость. Так как
$U(\lambda x_0)=\lambda^mU(x_0)<0$ при всех $\lambda>0$, нуль не минимум, и работает
лемма~\ref{lem:total-lyap}.
\end{proof}

Здесь $\dot J=\langle\dot x,x\rangle$ --- функция вида $\langle\dot x,v(x)\rangle$ с полем $v(x)=x$.
Сработали два свойства поля:
\begin{enumerate}[label=(\roman*)]
\item поле <<толкает вниз по потенциалу>>: $\langle v,\nabla U\rangle=mU<0$ на $\{U<0\}$;
\item поле <<растягивает>>: производная $Dv=\Id$ положительно определена.
\end{enumerate}
Для неоднородного $U$ радиальное поле ломается (пример~\ref{ex:radial-fails}), и вся идея
Паламодова --- найти поле со свойствами (i), (ii) для любого аналитического $U$ на плоскости.

\subsection{Критерий поля}

Для квадратной матрицы $M$ через $\sym M=\tfrac12(M+M^{\mathsf T})$ обозначим её симметричную
часть; $\sym M\succeq c$ значит $\langle M\xi,\xi\rangle\ge c|\xi|^2$ при всех $\xi$;
$\lVert M\rVert$ --- операторная норма. Поле $v$ с $\sym Dv\succeq0$ называют \term{монотонным}, с
$\sym Dv\succeq c>0$ --- \term{сильно монотонным}: тогда $\langle v(x)-v(y),x-y\rangle\ge c|x-y|^2$
для $x,y$ из выпуклого множества, где поле определено (интегрирование вдоль отрезка). Пример:
градиент выпуклой функции.

Ключевое вычисление. Пусть $v$ --- $C^1$-поле и $x(t)$ --- движение. Для
$F(t)=\langle\dot x(t),v(x(t))\rangle$
\begin{equation}\label{eq:Fdot}
\dot F=\langle\ddot x,v\rangle+\langle\dot x,Dv\,\dot x\rangle
=-\langle\nabla U,v\rangle+\langle\sym Dv\,\dot x,\dot x\rangle .
\end{equation}
Первый член неотрицателен, если поле толкает вниз; второй --- если поле монотонно.

\begin{theorem}[критерий поля]\label{thm:field}
\mk{Доказано; то же рассуждение проверено рецензентом: \texttt{notes/reviewer/2026-10-21-max-gluing.md},
лемма~1 из \texttt{notes/vlad-claude/theoremB-proof.md}; оформлено как \texttt{results/001}} Пусть
$U\in C^2$ около замкнутого шара $\bar B_\rho$, $U(0)=0$, $X=\{x\in B_\rho:U(x)<0\}$. Пусть на
открытом множестве, содержащем $X$, задано $C^1$-поле $A$ с $|A|\le M$ на $X$ и
\[
\sym DA(x)\succeq0,\qquad \langle A(x),\nabla U(x)\rangle\le\lambda\,U(x)\qquad(x\in X)
\]
при некотором $\lambda>0$. Тогда любое движение с энергией $E<0$, остающееся в $B_\rho$ при
$t\in[0,T]$, удовлетворяет $T\le2\beta M/(\lambda|E|)$, где $\beta=\sqrt{2(E-\min_{\bar B_\rho}U)}$.
В частности, $0$ тотально неустойчива.
\end{theorem}
\begin{proof}
На движении $U(x(t))\le E<0$, т.\,е. $x(t)\in X$, и $|\dot x|\le\beta$. По \eqref{eq:Fdot}
$\dot F\ge-\lambda U(x)\ge\lambda|E|$, а $|F|\le\beta M$. Значит, $2\beta M\ge F(T)-F(0)\ge\lambda|E|T$.
\end{proof}

\paragraph{Схема Паламодова 1977.} Доказательство в [P77] устроено так \mk{Литература: [P77], с.~42--45,
сверено по тексту; так же излагает Negrini 1995, §2}. Основная лемма: в роге $X^-_\rho$ есть
\emph{непрерывное} поле $v$ с $\langle v,\nabla U\rangle\le0$, гладкое вне конечного числа кривых и с
$\langle Dv\,\xi,\xi\rangle\ge c|\xi|^2$ вне них. Вспомогательный факт: при малом $\rho$ в роге нет
критических точек (лемма~\ref{lem:critval}). Затем берётся поле $v-\sigma\nabla U$ с малым $\sigma>0$
и функция $\langle v-\sigma\nabla U,\dot x\rangle$. Построение поля --- раздел~\ref{sec:plane}. Ниже ---
полное доказательство того, что из основной леммы следует тотальная неустойчивость, причём для
непрерывных кусочно гладких полей.

\begin{proposition}\label{prop:palamodov-variant}
\mk{Доказано} Пусть $U$ вещественно-аналитична около $0$, $U(0)=0$, а $\rho$ взято из
леммы~\ref{lem:critval}; $X=\{|x|<\rho/2,\ U<0\}$. Пусть на открытом множестве $\mathcal O\supset X$
задано ограниченное непрерывное поле $v$ и конечное объединение $K$ компактных вложенных
$C^1$-дуг, причём
$v\in C^1(\mathcal O\setminus K)$, $\sym Dv\succeq c>0$ на $\mathcal O\setminus K$ и
$\langle v,\nabla U\rangle\le0$ на $X$. Тогда любое движение с энергией $E<0$, остающееся в шаре
$|x|<\rho/2$ при $t\in[0,T]$, удовлетворяет
\[
T\le\frac{2\beta(M+\sigma L)}{\sigma\,\delta(E)^2},\qquad M=\sup_X|v|,\ L=\sup_X|\nabla U|,\
\sigma=\frac{c}{\sup_X\lVert D^2U\rVert}
\]
(любое $\sigma>0$, если $D^2U\equiv0$), а $\beta$ и $\delta(E)$ --- как выше. В частности, $0$
тотально неустойчива.
\end{proposition}

Доказательство опирается на две леммы о непрерывных полях.

\begin{lemma}\label{lem:monotone}
\mk{Доказано} В условиях предложения~\ref{prop:palamodov-variant} для любого выпуклого
$C\subset\mathcal O$ и $x,y\in C$ выполнено $\langle v(x)-v(y),x-y\rangle\ge c|x-y|^2$.
\end{lemma}
\begin{proof}
Пусть отрезок $[y,x]$ пересекает $K$ в конечном числе точек. Функция
$f(s)=\langle v(y+s(x-y)),x-y\rangle$ непрерывна на $[0,1]$ и дифференцируема вне конечного числа
точек, где $f'(s)=\langle Dv\,(x-y),x-y\rangle\ge c|x-y|^2$. По теореме Лагранжа на каждом
промежутке между этими точками и по непрерывности $f(1)-f(0)\ge c|x-y|^2$, что и требуется.
Общий случай: зафиксируем направление $e=(x-y)/|x-y|$ и сдвигаем отрезок параллельно. Для
$C^1$-дуги $\gamma$ проекция $\pi\circ\gamma$ на прямую, ортогональную $e$, --- $C^1$-функция
одной переменной; по теореме Сарда почти каждое её значение регулярно, и прообраз такого
значения на компактной дуге конечен. Значит, сколь угодно близкие параллельные отрезки (они лежат в $\mathcal O$: компактный отрезок
$[y,x]$ имеет окрестность внутри открытого $\mathcal O$) пересекают $K$
конечно (при $n\ge3$ почти все такие сдвиги вообще не пересекают $K$), и неравенство для
$[y,x]$ получается предельным переходом по непрерывности $v$.
\end{proof}

\begin{lemma}\label{lem:dini}
\mk{Доказано; рассуждение предложено рецензентом} Нижние производные Дини функции $F$ ---
это $\underline D^{\pm}F(t)=\liminf_{h\to0\pm}(F(t+h)-F(t))/h$. Пусть $v$ --- непрерывное поле на открытом
$\mathcal O$, удовлетворяющее заключению леммы~\ref{lem:monotone} на каждом шаре
$B\subset\mathcal O$ с константой $c\ge0$, и $x(t)$ --- движение, лежащее в $\mathcal O$ при
$t\in[0,T]$. Тогда для $F(t)=\langle\dot x(t),v(x(t))\rangle$ обе нижние производные Дини в каждой
точке не меньше $-\langle\nabla U(x),v(x)\rangle+c|\dot x|^2$, и, если эта величина $\ge\alpha$ при
всех $t$, то $F(T)-F(0)\ge\alpha T$.
\end{lemma}
\begin{proof}
Положим $\Delta x=x(t+h)-x(t)$, $\Delta v=v(x(t+h))-v(x(t))$. Тогда
\[
\frac{F(t+h)-F(t)}h=\Big\langle\frac{\dot x(t+h)-\dot x(t)}h,\,v(x(t+h))\Big\rangle
+\frac{\langle\Delta x,\Delta v\rangle}{h^2}+\Big\langle\dot x(t)-\frac{\Delta x}h,\,\frac{\Delta v}h\Big\rangle .
\]
Первое слагаемое стремится к $\langle\ddot x,v\rangle=-\langle\nabla U,v\rangle$. Второе при малом
$|h|$ (обе точки в малом шаре около $x(t)$) не меньше $c|\Delta x/h|^2\to c|\dot x|^2$ --- при любом
знаке $h$. Модуль третьего не больше $\tfrac12\sup|\ddot x|\cdot|\Delta v|\to0$ по непрерывности $v$.
Отсюда первое утверждение. Второе --- стандартный факт: непрерывная функция $G$ с нижней правой
производной Дини $>0$ всюду на $[0,T)$ не убывает; применяем его к $G=F-(\alpha-\eta)t$ и
устремляем $\eta\to0$.
\end{proof}

\begin{proof}[Доказательство предложения~\ref{prop:palamodov-variant}]
Положим $v_\sigma=v-\sigma\nabla U$. На $\mathcal O\setminus K$ (вблизи $X$)
$\sym Dv_\sigma=\sym Dv-\sigma D^2U\succeq c-\sigma\lVert D^2U\rVert\succeq0$, так что
лемма~\ref{lem:monotone} верна для $v_\sigma$ с константой $0$. Кроме того,
$\langle v_\sigma,\nabla U\rangle\le-\sigma|\nabla U|^2$. На движении с энергией $E<0$ точка $x(t)$
лежит в $\{|x|<\rho/2,\ U\le E\}$, где $|\nabla U|\ge\delta(E)>0$ (лемма~\ref{lem:critval}). По
лемме~\ref{lem:dini} для $F=\langle\dot x,v_\sigma(x)\rangle$ имеем
$F(T)-F(0)\ge\sigma\delta(E)^2T$, а $|F|\le\beta(M+\sigma L)$.
\end{proof}

\begin{remark}\label{rem:where-analytic}
Где здесь аналитичность: только в лемме~\ref{lem:critval} --- на подуровне $\{U\le E\}$ около
нуля нет критических точек. Вся остальная трудность --- \emph{построить} поле $v$. Если для
некоторого гладкого потенциала тотальная неустойчивость неверна (по Burgos--Paternain,
arXiv:2208.01139, таков пример Уреньи), то для него такого поля нет; значит, построение поля в
общем случае обязано использовать аналитичность ещё раз \mk{Гипотеза в той мере, в какой
опирается на вывод Burgos--Paternain о примере Уреньи}. Свобода склеивать поля непрерывно, без сглаживания,
существенна: при сглаживании разбиением единицы возникает член $(v_1-v_2)\otimes\nabla\chi$,
портящий монотонность. На склейке такого рода споткнулось отозванное доказательство Burgos 2021
для произвольного $n$ \mk{Литература: arXiv:2108.05829v4, лемма 2.2 --- поле $V=\sum f_pV_p$
склеено разбиением единицы; комментарий автора к v5: <<the second term of the last equation in
Lemma 2.2 does not have the desired asymptotic behavior>>}.
\end{remark}

\subsection{Почему радиальное поле не годится}

\begin{example}\label{ex:radial-fails}
\mk{Доказано, прямое вычисление} Для $U=(y-x^2)^2-x^6$ (пример~\ref{ex:horns}.4) положим
$s=y-x^2$. Для радиального поля $v=(x,y)$:
\[
\langle v,\nabla U\rangle=x(-4xs-6x^5)+2ys=2s^2-2x^2s-6x^6 .
\]
В нижней половине рога, $s=-\mu x^3$ с $0<\mu<1$, это равно $2\mu x^5+(2\mu^2-6)x^6>0$ при малых
$x>0$: радиальное поле толкает <<вверх по потенциалу>>, поперёк искривлённого рога
(рис.~\ref{fig:horn}). Поле $v=(x,2y)$, касающееся парабол $y=cx^2$, годится:
\[
\langle v,\nabla U\rangle=-4x^2s-6x^6+4ys=4s^2-6x^6=4U-2x^6\le4U,\qquad Dv=\operatorname{diag}(1,2),
\]
и по теореме~\ref{thm:field} равновесие тотально неустойчиво.
\end{example}

Поле $v=(x,a\,y)$ называют \term{взвешенным полем Эйлера} с весом $a$; его поток
$x\mapsto e^tx$, $y\mapsto e^{at}y$ сохраняет кривые $y=cx^a$. Если в рогу $U$ <<почти
квазиоднороден>> относительно весов $(1,a)$ (т.\,е. $U(\lambda x,\lambda^ay)\approx\lambda^dU(x,y)$),
такое поле работает. Но разные части одного рога могут требовать разных весов.

\begin{example}\label{ex:two-weights}
\mk{Доказано, прямое вычисление} $U=fg$, $f=y-x^2$, $g=y-x^3$; рог $\{x^3<y<x^2\}$, $x>0$. Для
$v=(x,a\,y)$ имеем $\langle v,\nabla f\rangle=af+(a-2)x^2$, $\langle v,\nabla g\rangle=ag+(a-3)x^3$,
откуда
\[
\langle v,\nabla U\rangle=2afg+(a-2)x^2g+(a-3)x^3f .
\]
На верхнем крае рога ($f=0$) это $x^4(a-2)(1-x)$, на нижнем ($g=0$) --- $x^5(3-a)(1-x)$. Там $U=0$,
так что при любом $\lambda\in\R$ условие $\langle v,\nabla U\rangle\le\lambda U$ требует $a\le2$ на
верхнем крае и $a\ge3$ на нижнем; по непрерывности при любом $a$ оно нарушается и в точках рога,
близких к одному из краёв. Ни одно поле $(x,a\,y)$ не годится: у верхнего края рог похож на кривую
веса $2$, у нижнего --- веса $3$, а посередине $U\approx-x^2y$.
\end{example}

Как Паламодов справляется с этим на плоскости --- содержание раздела~\ref{sec:plane}.

%=====================================================================
\section{Конструкция Паламодова на плоскости}\label{sec:plane}

В этом разделе изложено доказательство из [P77] по тексту статьи (русский оригинал,
\emph{Функц. анализ и его прил.} 11:4, с.~42--55). Обозначения приведены к нашим: потенциал $U$
(у Паламодова $u$), градиент $\nabla U$ (у Паламодова $u'$).

\subsection{Формулировка и план}

\begin{theorem}[Паламодов 1977, теорема 1]\label{thm:P77}
\mk{Литература: [P77], с.~42; сверено по тексту} Пусть $n=2$ и точка $x=0$ не является даже
нестрогим минимумом $U$, т.\,е. множество $\{U<U(0)\}$ непусто и его замыкание содержит нуль. Тогда равновесие $x=0$ неустойчиво, если выполнено одно из условий:
\begin{enumerate}[label=\Alph*)]
\item $U$ раскладывается в сходящийся степенной ряд в окрестности нуля;
\item $U$ бесконечно дифференцируема около нуля, и критическая точка $0$ конечнократна: в
кольце формальных степенных рядов идеал, порождённый рядами Маклорена функций
$\partial U/\partial x_1$, $\partial U/\partial x_2$, имеет конечную коразмерность.
\end{enumerate}
\end{theorem}

Изолированность в случае~A) не предполагается. В той же статье есть теорема~2 для любого $n$
(раздел~\ref{subsec:P77-rest}).

Пусть $B_\rho=\{|x|<\rho\}$, а $X^-_\rho$ --- связная компонента множества $\{U<0\}\cap B_\rho$,
замыкание которой содержит нуль (у нас --- <<рог>>; Паламодов формулирует чуть иначе, но в
доказательствах пользуется именно этим). Рогов конечное число (раздел~\ref{subsec:puiseux}).

\begin{lemma}[основная лемма]\label{lem:P77main}
\mk{Литература: [P77], с.~42--43} При некотором $\rho>0$ в $X^-_\rho$ можно задать
непрерывное векторное поле $v$ такое, что
\[
\text{(2)}\ \ \langle v,\nabla U\rangle\le0;\qquad\text{(3)}\ \ \langle Dv\,\xi,\xi\rangle\ge c|\xi|^2,\ c>0,
\]
причём $v$ имеет непрерывные первые производные всюду в $X^-_\rho$, кроме конечного числа
бесконечно дифференцируемых кривых, а (3) выполнено вне этих кривых.
\end{lemma}

\emph{Как из основной леммы следует теорема.} Это предложение~\ref{prop:palamodov-variant},
применённое к одному рогу: его доказательство проходит дословно с $\mathcal O=X^-_\rho$, потому что
движение с энергией $E<0$, пока оно в шаре, не покидает своей компоненты множества $\{U\le E\}$.
Для неустойчивости по Ляпунову хватает одного рога; для тотальной неустойчивости поле строится в
каждом из конечного числа рогов. Схема: поле $w=v-\sigma\nabla U$, функция
$l(t)=\langle w(x(t)),\dot x(t)\rangle$, $\dot l\ge\theta=\sigma\,\delta(E)^2>0$ на подуровне
$\{U\le E\}$ ($\delta(E)$ --- из леммы~\ref{lem:critval}), откуда время пребывания в $B_\rho$ не больше
$2\sqrt{2M_0}\,M'/\theta$, $M_0=\max_{|x|\le\rho}|U|$, $M'=\max_{|x|\le\rho}|w|$. Точки, где $v$ не
гладко, Паламодов обходит через меру: если бы траекторий, проводящих на исключительных кривых
время положительной меры, было несчётно много, среди них нашлись бы две, пересекающиеся по
подмножеству этих кривых положительной меры, и они совпали бы; для остальных траекторий оценка
получается суммированием и предельным переходом. Наша лемма~\ref{lem:dini} даёт то же короче. Остаётся
построить поле.

\subsection{Шаг 1. В роге нет критических точек}

\mk{Литература: [P77], с.~45; это лемма~\ref{lem:critval} при $n=2$} Пусть $U(0)=0$, $Z$ ---
множество нулей $U$, $\mathrm{Crit}$ --- множество её критических точек в $B_\rho$. У них нет внутренних точек
(иначе $U\equiv0$), поэтому это аналитические множества размерности $0$ или $1$. Уменьшив
$\rho$, можно считать, что каждое из них --- нуль и конечное число \term{элементарных кривых}
(образов аналитических отображений $(0,1)\to\R^2$, которые продолжаются аналитически на
окрестность отрезка $[0,1]$), выходящих из нуля. На кривой $l\subset\mathrm{Crit}$ функция $U$
постоянна, а по непрерывности равна $U(0)=0$. Значит, $\mathrm{Crit}\subset Z$ и $\nabla U\ne0$ в $X^-_\rho$.

\subsection{Шаг 2. Угол между $-x$ и $\nabla U$}

Пусть $\varphi(x)\in(-\pi,\pi]$ --- ориентированный угол от вектора $-x$ к вектору $\nabla U(x)$.
В полярных координатах $x=(r\cos\alpha,r\sin\alpha)$ пишем $\varphi=\varphi(\alpha,r)$.
Условие (2) для поля, повёрнутого относительно радиального на угол $\Delta$, --- это
$|\Delta-\varphi|\le\pi/2$; для самого радиального поля --- $|\varphi|\le\pi/2$: градиент не
смотрит наружу.

\begin{lemma}[угловая лемма; лемма 3 в P77]\label{lem:P77angle}
\mk{Литература: [P77], с.~45--46} При достаточно малом $\rho$ функция $\varphi$ не принимает в
$X^-_\rho$ значений $\pm\pi$, и
\begin{equation}\label{eq:angle5}
\sup\bigl\{|\varphi(x)|-\tfrac\pi2:\ x\in X^-_\rho,\ |x|=r\bigr\}\le c\,r^d
\end{equation}
с некоторыми $d>0$, $c\ge0$.
\end{lemma}

\emph{Доказательство качественной части} \mk{Доказано; по [P77], с.~45--46}: если $x_k\to0$,
$x_k\in X^-_\rho$ и $\varphi(x_k)\to\psi$, то $|\psi|\le\pi/2$. Пусть, напротив, $|\psi|>\pi/2$, и
$c=-\tfrac12\cos\psi>0$. При больших $k$ точки $x_k$ лежат в полуаналитическом множестве
$S=\{x\in B_\rho:\ U(x)<0,\ \langle x,\nabla U\rangle>0,\ \langle x,\nabla U\rangle^2>c^2|x|^2|\nabla U|^2\}$
(в нём угол между $x$ и $\nabla U$ острый с косинусом больше $c$). По лемме о выборе кривой есть
аналитическая дуга $f\colon[0,\eps)\to\R^2$, $f(0)=0$, $f(t)\in S$ при $t>0$; $f(t)=at^k+O(t^{k+1})$,
$a\ne0$, поэтому $f'(t)/|f'(t)|$ и $f(t)/|f(t)|$ стремятся к $a/|a|$. Значит,
$\frac{d}{dt}U(f(t))=\langle\nabla U,f'\rangle=|\nabla U||f'|\bigl(\cos\angle(\nabla U,f)+o(1)\bigr)>0$ при малых $t$:
$U$ возрастает вдоль дуги от нуля, и так как $U(f(0))=0$, на дуге $U>0$ --- противоречие с $f(t)\in S$.
(У Паламодова то же рассуждение проведено через полуаналитичность графика $\varphi$; на малых дугах
$S$ лежит в одном роге, так как $U<0$ на дуге.)

Количественная оценка \eqref{eq:angle5} \mk{по [P77], конспективно} получается так. Образ
относительно компактного полуаналитического множества при аналитическом отображении Паламодов
называет \term{конструктивным} множеством (сейчас говорят: \term{субаналитическое}); в размерности 2
это то же, что полуаналитическое [P77, свойство ($\theta$)]. Образ графика $\varphi$ при отображении
$(x,\alpha)\mapsto(|x|,|\alpha|)$ полуаналитичен, его верхняя граница при малых $r$ --- элементарная
кривая $\alpha=\alpha(r)$ с разложением Пюизё и $\alpha(0)\le\pi/2$; если $\alpha(0)=\pi/2$, то
$\alpha(r)-\pi/2=a\,r^{i/p}+\dots$, откуда \eqref{eq:angle5} с $d=i/p$.

Смысл: в роге $\partial_rU=|\nabla U|\cos\angle(x,\nabla U)=-|\nabla U|\cos\varphi\le\sin(cr^d)\,|\nabla U|$:
там, где потенциал растёт наружу (а это может быть и половина рога, см. пример~\ref{ex:radial-fails}),
градиент почти перпендикулярен радиусу.

\mk{Эскиз} Качественная часть переносится на любое $n$. Ориентированного угла в $\R^n$ нет; берём
неориентированный угол $\psi(x)\in[0,\pi]$ между $-x$ и $\nabla U$ и то же рассуждение с полуаналитическим
множеством $\{U<0,\ \langle x,\nabla U\rangle>0,\ \langle x,\nabla U\rangle^2>c^2|x|^2|\nabla U|^2\}$.

\subsection{Шаг 3. Повёрнутое радиальное поле}

Поле Паламодова задаётся в полярных координатах:
\begin{equation}\label{eq:Pfield}
|v(x)|=|x|,\qquad \arg v(x)=\arg x+\Delta(\arg x,|x|),
\end{equation}
т.\,е. это радиальное поле $x$, повёрнутое в каждой точке на угол $\Delta$. Ниже $\Delta'_\alpha$ и
$\Delta'_r$ --- частные производные по $\alpha$ и по $r$.

\begin{proposition}\label{prop:rotated}
\mk{Доказано; это вычисление из [P77], с.~52--53} Пусть $\Delta$ --- непрерывная функция на
$X^-_\rho$, гладкая вне конечного числа кривых, причём
\begin{enumerate}[label=(\Roman*)]
\item $\varphi-\tfrac\pi2\le\Delta\le\varphi+\tfrac\pi2$;
\item[(III)] $\lambda_H(r)\le\Delta'_\alpha\le\Lambda$, где $\Lambda>0$ --- константа, а $\lambda_H(r)\le0$,
$\lambda_H(r)\to0$ при $r\to0$;
\item[(IV)] $|r\Delta'_r|\to0$ и $\Delta\to0$ при $r\to0$ равномерно по $\alpha$.
\end{enumerate}
Тогда поле \eqref{eq:Pfield} удовлетворяет (2), а при малом $\rho$ --- и (3) с $c=\tfrac12$.
\end{proposition}
\begin{proof}
(2): по (I) угол между $v$ и $\nabla U$ не меньше $\pi/2$. (3) проверяем вне исключительных кривых,
где $\Delta$ гладкая (поэтому условие (II) статьи о гладкости здесь входит в предположения). Минимум
формы $\langle Dv\,\xi,\xi\rangle$ на единичной окружности не зависит от выбора ортонормированного
базиса, поэтому считаем в репере $(e_r,e_\alpha)$. Из $v=r(\cos(\alpha+\Delta),\sin(\alpha+\Delta))$ получаем
\[
Dv=\begin{pmatrix}\cos\Delta-r\Delta'_r\sin\Delta & -(1+\Delta'_\alpha)\sin\Delta\\
\sin\Delta+r\Delta'_r\cos\Delta & (1+\Delta'_\alpha)\cos\Delta\end{pmatrix}.
\]
Для матрицы $(a_{ij})$ минимум формы $\sum a_{ij}\xi_i\xi_j$ на единичной окружности равен
$\mu=\tfrac12\bigl(a_{11}+a_{22}-\sqrt{\sum a_{ij}^2-2\det(a_{ij})}\bigr)$ (лемма 8 в [P77]; это меньшее
собственное значение симметричной части). Подстановка даёт
\[
\mu=\cos\Delta-\tfrac12r\Delta'_r\sin\Delta+\tfrac12\Delta'_\alpha\cos\Delta-\tfrac12\sqrt{(r\Delta'_r)^2+(\Delta'_\alpha)^2}
\ \ge\ \cos\Delta-|r\Delta'_r|-\tfrac12(1-\cos\Delta)\Delta'_\alpha+\min(\Delta'_\alpha,0),
\]
По (IV) первые два слагаемые дают $1+o(1)$; третье по модулю не больше
$\tfrac12(1-\cos\Delta)\max(\Lambda,|\lambda_H|)=o(1)$; четвёртое не меньше $\lambda_H(r)=o(1)$ (здесь использовано $\tfrac12(q-|q|)=\min(q,0)$).
Значит, $\mu\ge\tfrac12$ при малом $r$.
\end{proof}

Заметим асимметрию: из точной формулы при $\Delta\to0$, $r\Delta'_r\to0$ и ограниченной
$\Delta'_\alpha$ получается $\mu=1+\min(\Delta'_\alpha,0)+o(1)$. Положительная производная
$\Delta'_\alpha$ может быть не малой, но должна быть ограничена константой $\Lambda$ (она входит с
множителем $1-\cos\Delta\to0$); отрицательная должна быть не меньше $-1+\kappa$ с $\kappa>0$.
Условие $\lambda_H(r)\to0$ --- выбор Паламодова в построении, с запасом. Поворот может нарастать
при обходе против часовой стрелки с ограниченной скоростью, а убывать --- только медленнее, чем
$-\Delta'_\alpha<1$.

\subsection{Шаг 4. Построение угла поворота}

Если при некотором $\rho$ в $X^-_\rho$ выполнено $|\varphi|\le\pi/2$, достаточно $v(x)=x$. Иначе
вводится минимально нужный поворот (у Паламодова $\delta$)
\[
\tau(\alpha,r)=\begin{cases}\varphi-\tfrac\pi2,&\varphi\ge\tfrac\pi2,\\ 0,&|\varphi|<\tfrac\pi2,\\
\varphi+\tfrac\pi2,&\varphi\le-\tfrac\pi2.\end{cases}
\]
Он удовлетворяет (I), а по угловой лемме $|\tau|\le cr^d\to0$. Но $\Delta=\tau$ брать нельзя:
поперёк узкого рога $\varphi$ меняется быстро, и $\tau$ может меняться по $\alpha$ слишком круто.
Поэтому $\tau$ заменяют огибающей с ограниченными наклонами. Возможность этого обеспечивает
следующая лемма о том, как меняется $\varphi$ вдоль дуги окружности.

\begin{lemma}[лемма 5 в P77]\label{lem:P77five}
\mk{Литература: [P77], с.~46--48; доказательство там опирается на полуаналитичность, лемму 4
($r\,d\alpha/dr=\operatorname{tg}$ угла между радиусом и касательной к кривой $\alpha=\alpha(r)$) и
геометрию двух близких кривых; здесь не воспроизводится} Существуют $\rho,\gamma_1,\eta_1>0$
такие, что для точек $(\alpha,r),(\beta,r)\in X^-_\rho$
\[
\varphi(\beta,r)-\varphi(\alpha,r)\ge-\pi\ \ \text{при}\ \alpha\le\beta\le\alpha+\gamma_1,\qquad
\varphi(\beta,r)-\varphi(\alpha,r)\le\pi+\eta_1(\beta-\alpha)\ \ \text{при}\ \alpha\le\beta.
\]
\end{lemma}

\emph{Огибающая.} Положим $\lambda_B=\eta_1$ (допустимый подъём) и $\lambda_H(r)=-2cr^d/\gamma_1$
(допустимый спуск; увеличив $\eta_1$, добиваются $cr^d\le\frac12\eta_1\gamma_1$). Точки, где
$\tau'_\alpha=\lambda_B$ или $\tau'_\alpha=\lambda_H(r)$, образуют конструктивное множество и при
малом $\rho$ заполняют конечное число элементарных кривых $\alpha=\alpha_i(r)$, $i=1,\dots,I$. Здесь
производная понимается в обобщённом смысле: в изломах $\tau$ (где $|\varphi|=\pi/2$) производной нет,
и Паламодов берёт замыкание множеств, где разностные отношения равны $\lambda+\eps$, $|\eps|\le1$,
пересечённое с $\{\beta=\alpha,\ \eps=0\}$; так в число кривых попадают и изломы, в которых
разностные отношения <<перескакивают>> через $\lambda$. Из точки $(\alpha_i(r),\tau(\alpha_i(r),r))$ проводят
прямую (по $\alpha$) с наклоном $\lambda_B$ или $\lambda_H$ до пересечения с нулём; получается функция
$\Delta_i$, линейная на этом отрезке и равная нулю вне его, с носителем длины не больше
$\frac12\gamma_1$. В точке $\alpha_i(r)$ функция $\Delta_i$ имеет скачок, но там $\Delta_i=\tau$, поэтому
итоговая $\Delta$ остаётся непрерывной. Наконец, $\Delta_0=\tau$ и
\[
\Delta=\max\{\Delta_i:\Delta_i>0\},\quad\text{если есть }\Delta_i>0;\qquad
\Delta=\min\{\Delta_i:\Delta_i<0\},\quad\text{если есть }\Delta_i<0;\qquad\Delta=0\ \text{иначе}.
\]
Подъём $\tau$ круче $\lambda_B$ заранее <<подготавливается>> пандусом наклона $\lambda_B$ слева, а
спуск круче $\lambda_H$ заменяется пологим спуском наклона $\lambda_H$ справа.

\begin{figure}[ht]
\centering
\begin{tikzpicture}[xscale=1.6,yscale=1.4]
\draw[->,gray] (-0.2,0)--(6.4,0) node[right]{$\alpha$};
\draw[thick] (0,0) -- (0.9,0) .. controls (1.05,0.05) and (1.15,0.75) .. (1.3,0.8) .. controls (1.45,0.85) and (1.55,0.1) .. (1.65,0) -- (2.6,0) .. controls (2.75,0.08) and (2.85,0.55) .. (3.0,0.55) .. controls (3.15,0.55) and (3.2,0.05) .. (3.3,0) -- (4.4,0) .. controls (4.5,-0.05) and (4.6,-0.55) .. (4.75,-0.55) .. controls (4.9,-0.55) and (5.0,-0.05) .. (5.1,0) -- (6.2,0);
\draw[dashed,thick,blue] (0.55,0) -- (1.18,0.62);
\draw[dashed,thick,blue] (1.42,0.78) -- (2.55,0.25) -- (2.62,0.22);
\draw[dashed,thick,blue] (2.62,0.22) -- (2.82,0.42);
\draw[dashed,thick,blue] (3.08,0.53) -- (4.75,-0.55);
\draw[dashed,thick,blue] (4.95,-0.42) -- (5.5,0);
\node at (1.3,1.05) {\small $\delta$};
\node[blue] at (2.2,0.75) {\small $\Delta$};
\end{tikzpicture}
\caption{Схема огибающей (как рис.~3 в [P77]): сплошная линия --- минимально нужный поворот
$\tau$ при фиксированном $r$, пунктир --- функция $\Delta$ там, где она отличается от $\delta$:
пологие подъёмы наклона не больше $\lambda_B$ и медленные спуски наклона $\lambda_H(r)\to0$.}
\label{fig:envelope}
\end{figure}

\begin{lemma}[леммы 6 и 7 в P77]\label{lem:P77seven}
\mk{Литература: [P77], с.~49--52; доказательство изложено конспективно} (а) В каждой точке все
ненулевые $\Delta_i$ одного знака, и $0<\Delta_i\le\frac\pi2+\varphi$ или $-\frac\pi2+\varphi\le\Delta_i<0$
(здесь используется лемма~\ref{lem:P77five}). (б) Функция $\Delta$ корректно определена,
непрерывна, полуаналитична и удовлетворяет условиям (I), (III), (IV) предложения~\ref{prop:rotated};
она дифференцируема всюду, кроме конечного числа элементарных кривых.
\end{lemma}

Самое тонкое --- оценка (IV) для радиальной производной. Паламодов сначала получает грубую
оценку $|\Delta'_r|\le Cr^A$: раскладывает $U=w\,u_1^{q_1}\cdots u_m^{q_m}$ на неприводимые
множители в кольце аналитических функций (каждая граничная кривая рога лежит в нулях одного
множителя) и оценивает производную по $r$ от направления $\nabla U/|\nabla U|$. Затем для
$\tilde\Delta=r^k\Delta$ с подходящим $k$ берёт полуаналитическую верхнюю границу
$z(r)\sim c\,r^\mu$ функции $|\tilde\Delta'_r|$ и показывает $\mu>k-1$: вдоль кривой, где
$|\tilde\Delta'_r|\ge\frac12z(r)$, интеграл $\int r^{-k}\tilde\Delta'_r\,dr$ ограничен, потому что
$\Delta$ ограничена, $\Delta'_\alpha$ ограничена по (III), $\int|\alpha'(r)|\,dr<\infty$ для кривой Пюизё, а
$\Delta=O(r^d)$ по угловой лемме.

\emph{Итог.} Поле \eqref{eq:Pfield} с этим $\Delta$ удовлетворяет (2) и (3)
(предложение~\ref{prop:rotated}); основная лемма в случае~A) доказана.

\subsection{Случай Б) и теорема 2}\label{subsec:P77-rest}

\emph{Случай Б)} \mk{Литература: [P77], с.~53--54; изложено кратко}. Из конечной кратности
(подготовительная теорема Мальгранжа) следует $|\nabla U(x)|\ge c|x|^b$, а по теореме Мазера о
достаточности струй гладкая замена $y=x+h(x)$ приводит $U$ к её многочлену Тейлора $\tilde u$
достаточно большой степени. Для $\tilde u$ поле $\tilde v$ даёт случай~A), затем берётся
$w=\tilde v-\sigma|y|^{2-m}\nabla\tilde u$ и поле $v(x)=y'(x)^{\mathsf T}w(y(x))$ в исходных
координатах; поправки от замены оцениваются через $|h'|=o(|y|^{1-m}|\nabla\tilde u|)$.

\emph{Теорема 2} \mk{Литература: [P77], с.~54--55}: в $\R^n$ пусть $U$ имеет непрерывные
производные до порядка $m$ (аналитичность не нужна), её дифференциалы в нуле порядка меньше $m$
равны нулю, а $m$-й дифференциал $u_m$ (как многочлен) не имеет ненулевых вещественных критических
точек, и нуль не его минимум. Тогда равновесие неустойчиво; более того, при начальном смещении
$x(0)$ с $U(x(0))<0$ и начальной скорости $|\dot x(0)|\le c\sqrt{-U(x(0))}$ с малым $c<1$ (условие (18)
статьи) траектория покидает шар $\{|x|\le\rho\}$ за время $T\sim\ln(1/|x(0)|)$ при $m=2$ и
$T\sim|x(0)|^{-(m-2)/2}$ при $m>2$; здесь $T\sim a$ значит $ca\le T\le Ca$. Поле здесь явное: $v=x-\sigma|x|^{2-m}\nabla U$. Это
предшественник теоремы 3 обзора [P20] и градиентного поля в \texttt{results/001}.

\subsection{Где работает аналитичность и что мешает при $n\ge3$}

Аналитичность используется трижды: (1) в роге нет критических точек (шаг 1); (2) угловая лемма
и её оценка $cr^d$ (лемма о выборе кривой и разложения Пюизё); (3) конечность числа кривых
$\alpha_i(r)$ и оценка (IV) (полуаналитичность всех построенных функций, факторизация $U$). В
$C^\infty$-примере Пенлеве (у Паламодова --- радиальный, $U=\exp(-|x|^{-1})\sin(|x|^{-1})$, множество
$\{U<0\}$ --- последовательность сферических слоёв) ломается уже начало: у $\{U<0\}$ нет ни одной
связной компоненты, замыкание которой содержит нуль, так что рог даже не определён.

\mk{Гипотеза главного агента} При $n\ge3$ первые два шага переносятся: доказательство угловой
леммы опирается только на лемму о выборе кривой. Ломается шаг 4. Угол поворота --- функция на
дуге окружности, а огибающая строится, двигаясь по дуге в одном направлении (наклоны
$\lambda_B$, $\lambda_H$ различают <<против>> и <<по>> часовой стрелке). В $\R^3$ сечение рога сферой ---
двумерная область, поворот радиального поля задаётся уже не одним числом, а аналога порядка на
дуге и леммы~\ref{lem:P77five} нет. Это кандидат на исследование на этапе 2.

%=====================================================================
\section{Что было дальше и почему не получается при $n\ge3$}\label{sec:after}

\paragraph{Другие доказательства при $n=2$.}
\begin{itemize}[leftmargin=*]
\item Талиаферро (1980) независимо доказал неустойчивость для аналитического $U$ на плоскости
\mk{Литература: S.\,D.~Taliaferro, J.~Differential Equations 35 (1980), 248--265 (несмотря на
название <<Stability…>>); по вторичным источникам}.
\item Козлов (1981) перенёс результат на произвольную кинетическую энергию, опираясь на работу
Паламодова \mk{Литература: В.\,В.~Козлов, УМН 36:1 (1981), 209--210; по вторичным источникам}.
\item Лалуа и Пайфер (1982) разобрали нестрогие минимумы на плоскости; вместе с предыдущими это
даёт полное обращение теоремы Лагранжа--Дирихле для аналитических $U$ при $n=2$
\mk{Литература: Burgos--Paternain, arXiv:2208.01139, введение}.
\item Брюнелла (1998) дал другое, топологическое доказательство (A) при $n=2$. Его теорема: если у
аналитического векторного поля с изолированной особой точкой есть неприводимое трёхмерное
инвариантное аналитическое множество, гладкое вне особой точки, с линком ненулевой эйлеровой
характеристики, то есть траектория, стремящаяся к особой точке. Следствие: если $0$ ---
изолированная критическая точка аналитического $U$ на плоскости, \emph{не минимум и не
максимум}, то есть асимптотическая траектория. При $n=2$ нулевой уровень энергии $\{H=0\}$
трёхмерен (\term{линк} --- пересечение множества с малой сферой вокруг особой точки). Случай
максимума (линк --- тор, $\chi=0$) Брюнелла лишь отмечает как элементарный (принцип Мопертюи и
теорема Хопфа--Ринова) \mk{Литература: M.~Brunella, Ann.\ Inst.\ Fourier 48 (1998),
1345--1357; сверено по тексту на Numdam}.
\end{itemize}

\paragraph{Попытки для всех $n$.} В 1992--1995 годах Паламодов анонсировал тотальную
неустойчивость для всех $n$ (через разрешение особенностей Хиронаки); полного доказательства
опубликовано не было \mk{Литература: Allaire--Rauch, ARMA 224 (2017), §1.7; Burgos,
arXiv:2108.05829v4, введение}. Оценки этой работы расходятся: Негрини (1995) и Burgos--Maderna--Paternain
(2021) принимали её как доказательство, Burgos (2021), Burgos--Paternain (2022) и сам Паламодов в
обзоре 2020 года --- нет \mk{Литература: см. \texttt{problem/KNOWN\_RESULTS.md}, §5}. Препринт
Burgos (2021) с доказательством для всех $n$ отозван автором из-за ошибки в лемме~2.2 о склейке
локальных полей \mk{Литература: arXiv:2108.05829, комментарий к v5}. В работе Burgos (2025)
тотальная неустойчивость доказана, в частности, когда ограничение младшей формы на единичную
сферу --- \term{функция Морса} (все её критические точки невырождены) и $0$ --- её
\term{регулярное значение} (в точках нулевого уровня дифференциал не равен нулю); случай (A)
назван открытым \mk{Литература: arXiv:2506.19135, предложение 1.2 и §1.2}.

\paragraph{Что ломается при $n\ge3$.} На плоскости множество $\{U<0\}$ около нуля --- конечное
число рогов, каждый ограничен двумя кривыми с явными разложениями Пюизё. В $\R^3$ это области,
ограниченные поверхностями; их сечения сферой $|x|=r$ --- области на двумерной сфере с
нетривиальной топологией, а явной параметризации нет. Нужны \term{раздутия} (замены координат,
заменяющие точку или кривую на сферу направлений или трубку; из них составляется разрешение
особенностей), а при переходе через раздутие монотонность поля ($\sym Dv\succeq0$) не
сохраняется, и склейка локальных полей вносит член без нужной малости. Это и есть открытая
часть задачи.

%=====================================================================
\appendix
\section{Словарь терминов}\label{app:glossary}

\begin{description}[leftmargin=0pt,style=nextline]
\item[Асимптотическая траектория] движение $x(t)\not\equiv0$, стремящееся вместе со скоростью к
равновесию при $t\to+\infty$ или $t\to-\infty$; её существование влечёт неустойчивость
(лемма~\ref{lem:asympt}).
\item[Взвешенное поле Эйлера] $v=(a_1x_1,\dots,a_nx_n)$; для квазиоднородного $U$ весов $(a_i)$ и
степени $d$ выполнено $\langle v,\nabla U\rangle=dU$.
\item[Вириал] член $\langle x,\nabla U\rangle$ в тождестве Лагранжа--Якоби.
\item[Квазиоднородная функция] $U(\lambda^{a_1}x_1,\dots,\lambda^{a_n}x_n)=\lambda^dU(x)$ при всех
$\lambda>0$.
\item[Критическая точка] точка, где $\nabla U=0$; изолированная --- если других рядом нет.
\item[Кусочно гладкое поле] непрерывное поле, гладкое вне конечного числа кривых (в
предложении~\ref{prop:palamodov-variant} --- вложенных $C^1$-дуг).
\item[Конструктивное (субаналитическое) множество] образ относительно компактного
полуаналитического множества при аналитическом отображении; в размерности 2 --- то же, что
полуаналитическое.
\item[Лемма о выборе кривой] если точка лежит в замыкании полуаналитического множества, к ней
ведёт аналитическая дуга из этого множества (теорема~\ref{thm:curve}).
\item[Линк] пересечение множества с малой сферой вокруг особой точки.
\item[Младшая форма] первая ненулевая однородная часть ряда Тейлора.
\item[Момент инерции] $|x|^2$ (у нас $J=\tfrac12|x|^2$).
\item[Монотонное поле] поле $v$ с $\sym Dv\succeq0$; сильно монотонное --- $\succeq c>0$.
\item[Натуральная механическая система] энергия <<кинетическая (квадратичная по скоростям) плюс
потенциальная (зависит только от положения)>>.
\item[Неустойчивость по Ляпунову] отрицание устойчивости (определение~\ref{def:stab}).
\item[Операторная норма] $\lVert M\rVert=\max_{|\xi|=1}|M\xi|$.
\item[Производные Дини] нижние: $\underline D^{\pm}F(t)=\liminf_{h\to0\pm}(F(t+h)-F(t))/h$.
\item[Полуаналитическое множество] конечное объединение множеств, заданных аналитическими
равенствами и строгими неравенствами.
\item[Полуветвь] одна из конечного числа дуг, на которые распадается $\{U=0\}$ около нуля на
плоскости (теорема~\ref{thm:puiseux}).
\item[Потенциал, потенциальная энергия] функция $U$; сила равна $-\nabla U$. Силовая функция ---
$-U$.
\item[Раздутие] замена координат, заменяющая точку (или кривую) на множество направлений в ней;
из раздутий составляется разрешение особенностей.
\item[Равновесие (положение равновесия)] постоянное решение; точка $(x_0,0)$ фазового
пространства с $\nabla U(x_0)=0$.
\item[Регулярное значение] значение функции, в прообразе которого дифференциал нигде не равен
нулю.
\item[Рог] связная компонента множества $\{U<0\}$ в малой проколотой окрестности нуля.
\item[Ряд Пюизё] сходящийся ряд по дробным степеням $x^{1/N}$.
\item[Строгий, нестрогий минимум] строгий: $U>U(0)$ в проколотой окрестности; нестрогий:
$U\ge U(0)$, но не строгий.
\item[Тотальная неустойчивость] все движения с энергией ниже $U(0)$ покидают фиксированную
окрестность за время, ограниченное через энергию (определение~\ref{def:total}).
\item[Устойчивость по Ляпунову] малые возмущения положения и скорости остаются малыми навсегда
(определение~\ref{def:stab}).
\item[Фазовое пространство] пространство пар (положение, скорость).
\item[Функция Морса] гладкая функция, все критические точки которой невырождены (гессиан
невырожден).
\item[Функция Четаева] функция на фазовом пространстве, положительная на области $D$ (компоненте
$\{V>0\}$), к границе которой примыкает равновесие, и растущая вдоль движений в $D$
(теорема~\ref{thm:chetaev}).
\item[$\sym M$] симметричная часть матрицы, $\tfrac12(M+M^{\mathsf T})$.
\end{description}

%=====================================================================
\begin{thebibliography}{99}
\bibitem[P77]{P77} В.\,П.~Паламодов, Об устойчивости равновесия в потенциальном поле,
\emph{Функц. анализ и его прил.} 11:4 (1977), 42--55; англ.: \emph{Funct. Anal. Appl.} 11
(1977), 277--289, doi:10.1007/BF01077142.
\bibitem[P95]{P95} V.\,P.~Palamodov, Stability of motion and algebraic geometry, in:
\emph{Dynamical Systems in Classical Mechanics} (V.\,V.~Kozlov, ed.), Amer. Math. Soc. Transl.
Ser.~2, vol.~168, 1995, 5--20.
\bibitem[P20]{P20} В.\,П.~Паламодов, Об обращении теоремы Лагранжа--Дирихле и неустойчивости
консервативных систем, \emph{УМН} 75:3 (2020), 107--122.
\bibitem[Neg95]{Neg95} P.~Negrini, On the inversion of Lagrange--Dirichlet theorem,
\emph{Resenhas IME-USP} 2:1 (1995), 83--114.
\bibitem[Ta80a]{Ta80a} S.\,D.~Taliaferro, An inversion of the Lagrange--Dirichlet stability theorem,
\emph{Arch. Rational Mech. Anal.} 73 (1980), 183--190.
\bibitem[Ta80b]{Ta80b} S.\,D.~Taliaferro, Stability for two dimensional analytic potentials,
\emph{J. Differential Equations} 35 (1980), 248--265.
\bibitem[Hag71]{Hag71} P.~Hagedorn, Die Umkehrung der Stabilitätssätze von Lagrange-Dirichlet und
Routh, \emph{Arch. Rational Mech. Anal.} 42 (1971), 281--316.
\bibitem[Koz81]{Koz81} В.\,В.~Козлов, Неустойчивость равновесия в потенциальном поле, \emph{УМН}
36:1 (1981), 209--210.
\bibitem[LP82]{LP82} M.~Laloy, K.~Peiffer, On the instability of equilibrium when the potential
has a non-strict local minimum, \emph{Arch. Rational Mech. Anal.} 78 (1982), 213--222.
\bibitem[Lal76]{Lal76} M.~Laloy, On equilibrium instability for conservative and partially
dissipative systems, \emph{Int. J. Non-Linear Mech.} 11 (1976), 295--301.
\bibitem[Ure20]{Ure20} A.\,J.~Ureña, To what extent are unstable the maxima of the potential?,
\emph{Ann. Mat. Pura Appl.} 199 (2020), 1763--1775.
\bibitem[BMP21]{BMP21} J.\,M.~Burgos, E.~Maderna, M.~Paternain, On the Lyapunov instability in
Newtonian dynamics, \emph{Nonlinearity} 34 (2021), 6719--6726; arXiv:1906.11962.
\bibitem[BP22]{BP22} J.\,M.~Burgos, M.~Paternain, On the Lagrange--Dirichlet converse in
dimension three, arXiv:2208.01139 (2022), препринт.
\bibitem[B21]{B21} J.\,M.~Burgos, A proof of the Palamodov's total instability conjecture,
arXiv:2108.05829 (2021; отозвано автором в 2022).
\bibitem[B25]{B25} J.\,M.~Burgos, McGehee blowup for Lagrangian systems and instability of
equilibria, \emph{J. Differential Equations} 454 (2026); arXiv:2506.19135.
\bibitem[Br98]{Br98} M.~Brunella, Instability of equilibria in dimension three, \emph{Ann. Inst.
Fourier} 48 (1998), 1345--1357.
\bibitem[AR17]{AR17} G.~Allaire, J.~Rauch, \emph{Arch. Rational Mech. Anal.} 224 (2017), 233--268.
\bibitem[EPS26]{EPS26} A.~Enciso, D.~Peralta-Salas, Stability and instability for particle
dynamics in magnetic fields, arXiv:2610.27893 (2026).
\bibitem[Mil68]{Mil68} J.~Milnor, \emph{Singular points of complex hypersurfaces}, Princeton, 1968.
\bibitem[Wall04]{Wall04} C.\,T.\,C.~Wall, \emph{Singular points of plane curves}, Cambridge, 2004.
\end{thebibliography}

\end{document}
